% lamb.tex with one property added: a dependency cycle a build must cut.
%
% PROPERTY — two lemmas whose proofs cite each other. Each proof opens with the
% plain "Proof." head, naming no result, so it binds to the lemma directly above
% it; each cites the other lemma in its body. A --no-inference build drafts each
% contained reference as "supported by", giving a supports edge each way. The
% cycle walk at the end of depends-attributed cuts one into a demoted record,
% and analyses-split then moves that record onto the analysis it splits from
% its lemma. Keep the references out of the proofs' opening runs: a reference
% there would bind the proof to the lemma it names and the cycle would not form.
%
% No bibliography travels with this file and it cites nothing: lamb.tex covers
% citations, and this fixture isolates the cycle.
\documentclass{article}
\usepackage{amsthm}
\newtheorem{lemma}{Lemma}
\title{Mary and Her Lamb: A Mutual Account}
\author{A. Nonymous}
\begin{document}
\maketitle
\begin{abstract}
We recount the case of Mary, who had a little lamb, and prove that each went
wherever the other did.
\end{abstract}

\section{Introduction}
Mary had a little lamb. Its fleece was white as snow.

\section{Results}
\begin{lemma}\label{lem:follows}
Everywhere that Mary went, the lamb was sure to go.
\end{lemma}
\begin{proof}
Let Mary go somewhere. Mary is never seen without the lamb, by
Lemma~\ref{lem:accompanied}, so the lamb is there too.
\end{proof}

\begin{lemma}\label{lem:accompanied}
Mary is never seen without the lamb.
\end{lemma}
\begin{proof}
Wherever Mary is seen, she went there. The lamb goes wherever Mary goes, by
Lemma~\ref{lem:follows}, so it is seen with her.
\end{proof}

\section{Conclusion}
It waited patiently about till Mary did appear.
\end{document}
